\documentclass[12pt,a4paper]{article}
\usepackage[T1]{fontenc}
\usepackage{amsmath,amssymb,mathrsfs,tikz,times,pifont}
\usepackage{enumitem}
\usepackage{multicol}
\usepackage{lmodern}
\usetikzlibrary{trees}
\newcommand\circitem[1]{%
\tikz[baseline=(char.base)]{
\node[circle,draw=gray, fill=red!55,
minimum size=1.2em,inner sep=0] (char) {#1};}}
\newcommand\boxitem[1]{%
\tikz[baseline=(char.base)]{
\node[fill=cyan,
minimum size=1.2em,inner sep=0] (char) {#1};}}
\setlist[enumerate,1]{label=\protect\circitem{\arabic*}}
\setlist[enumerate,2]{label=\protect\boxitem{\alph*}}
%%%::::::by chnini ameur :::::::%%%
\everymath{\displaystyle}
\usepackage[left=1cm,right=1cm,top=1cm,bottom=1.7cm]{geometry}
\usepackage[colorlinks=true, linkcolor=blue, urlcolor=blue, citecolor=blue]{hyperref}
\usepackage{array,multirow}
\usepackage[most]{tcolorbox}
\usepackage{varwidth}
\usepackage{float} %pour utiliser l'option [H] qui force l'image à apparaître exactement à l'endroit où elle est placée dans le code.
 \usetikzlibrary{calc}
\tcbuselibrary{skins,hooks}
\usetikzlibrary{patterns}
%%%::::::by chnini ameur :::::::%%%
\newtcolorbox{exa}[2][]{enhanced,breakable,before skip=2mm,after skip=5mm,
colback=yellow!20!white,colframe=black!20!blue,boxrule=0.5mm,
attach boxed title to top left ={xshift=0.6cm,yshift*=1mm-\tcboxedtitleheight},
fonttitle=\bfseries,
title={#2},#1,
% varwidth boxed title*=-3cm,
boxed title style={frame code={
\path[fill=tcbcolback!30!black]
([yshift=-1mm,xshift=-1mm]frame.north west)
arc[start angle=0,end angle=180,radius=1mm]
([yshift=-1mm,xshift=1mm]frame.north east)
arc[start angle=180,end angle=0,radius=1mm];
\path[left color=tcbcolback!60!black,right color = tcbcolback!60!black,
middle color = tcbcolback!80!black]
([xshift=-2mm]frame.north west) -- ([xshift=2mm]frame.north east)
[rounded corners=1mm]-- ([xshift=1mm,yshift=-1mm]frame.north east)
-- (frame.south east) -- (frame.south west)
-- ([xshift=-1mm,yshift=-1mm]frame.north west)
[sharp corners]-- cycle;
},interior engine=empty,
},interior style={top color=yellow!5}}
%%%%%%%%%%%%%%%%%%%%%%%
\usepackage{fancyhdr}
\usepackage{eso-pic}         % Pour ajouter des éléments en arrière-plan
% Commande pour ajouter du texte en arrière-plan
\usepackage{tkz-tab}
\AddToShipoutPicture{
    \AtTextCenter{%
        \makebox[0pt]{\rotatebox{80}{\textcolor[gray]{0.7}{\fontsize{5cm}{5cm}\selectfont PGB}}}
    }
}
\usepackage{lastpage}
\fancyhf{}
\pagestyle{fancy}
\renewcommand{\footrulewidth}{1pt}
\renewcommand{\headrulewidth}{0pt}
\renewcommand{\footruleskip}{10pt}
\fancyfoot[R]{
\color{blue}\ding{45}\ \textbf{2025}
}
\fancyfoot[L]{
\color{blue}\ding{45}\ \textbf{Prof:M. BA}
}
\cfoot{\bf
\thepage /
\pageref{LastPage}}
% Création du compteur pour les exercices
\newcounter{exercice}
\renewcommand{\theexercice}{\arabic{exercice}}  % Définit l'affichage du compteur en chiffres arabes

% Définir la commande \exo
\newcommand{\exo}{\refstepcounter{exercice}\textbf{Exercice \theexercice} }
\begin{document}
\renewcommand{\arraystretch}{1.5}
\renewcommand{\arrayrulewidth}{1.2pt}
\begin{tikzpicture}[overlay,remember picture]
    \node[draw=blue,line width=1.2pt,fill=purple,text=blue,inner sep=3mm,rounded corners,pattern=dots]at ([yshift=-2.5cm]current page.north) {\begingroup\setlength{\fboxsep}{0pt}\colorbox{white}{\begin{tabular}{|*1{>{\centering \arraybackslash}p{0.28\textwidth}} |*2{>{\centering \arraybackslash}p{0.2\textwidth}|} *1{>{\centering \arraybackslash}p{0.19\textwidth}|} }
                \hline
                \multicolumn{3}{|c|}{$\diamond$$\diamond$$\diamond$\ \textbf{Lycée de Dindéfélo}\ $\diamond$$\diamond$$\diamond$ } & \textbf{A.S. : 2025/2026}                                                                     \\ \hline
                \textbf{Matière: Mathématiques}                                                                                    & \textbf{Niveau : 1}\textbf{$^{er}$S2} & \textbf{Date: ../02/2026} & \textbf{Durée : 4 heures} \\ \hline
                \multicolumn{4}{|c|}{\parbox[c]{10cm}{\begin{center}
                                                                  \textbf{{\Large\sffamily Composition n$ ^{\circ} $ 1 Du 1$ ^\text{\bf er} $ Semestre}}
                                                              \end{center}}}                                                                                                                               \\ \hline
            \end{tabular}}\endgroup};
\end{tikzpicture}
\vspace{3cm}
\section*{\underline{ Exercice 1 :} (4,75 pts) \small\itshape Barycentres et lieux géométriques dans le parallélogramme}

Soit $ABCD$ un parallélogramme et soit $G$ le barycentre du système $\{(A,3);(B,2);(C,3);(D;2)\}$.

\begin{enumerate}
    \item Construire les barycentres $E$ et $F$ des systèmes respectifs $\{(A,3);(B,2)\}$ et $\{(C,3);(D;2)\}$ \hfill \textbf{1pt}
    
    \item Démontrer que $G$ est le milieu de $[EF]$, puis construire le point $G$. \hfill \textbf{0,5pt+0,25pt}
    
    \item Soit $I$ le centre de gravité de $ABC$ et $J$ le symétrique de $B$ par rapport à $D$
    \begin{enumerate}
        \item Construire les points $I$ et $J$. \hfill \textbf{0,5pt}
        \item Montrer que $J = \text{bar}\{(B;-1);(D;2)\}$. \hfill \textbf{0,5pt}
        \item Montrer que $I, J$ et $G$ sont alignés. \hfill \textbf{0,5pt}
    \end{enumerate}
    
    \item Déterminer et construire l'ensemble des points $M$ du plan tels que :\hfill \textbf{1,5pt}
    \begin{multicols}{2}
\setlength{\columnseprule}{0.9mm} % La largeur de la ligne verticale entre les colonnes
    \begin{enumerate}
        \item $\left\|3\overrightarrow{MA} + 2\overrightarrow{MB}\right\| = \left\|3\overrightarrow{MC} + 2\overrightarrow{MD}\right\|$.
        \item $\left\|3\overrightarrow{MA} + 2\overrightarrow{MB} + 3\overrightarrow{MC} + 2\overrightarrow{MD}\right\| = 20$.
    \end{enumerate}
    \end{multicols}
\end{enumerate}

\section*{\underline{ Exercice 2 :} (4,75 pts) \small\itshape Fonctions réelles : domaines, compositions et valeurs absolues}
\begin{enumerate}
\item Déterminer le domaine de définition $D_f$ de la fonction $f$ définie ci-dessous \textbf{0,5pt+0,5pt+0,5pt+1pt}
\begin{multicols}{2}
\setlength{\columnseprule}{0.5mm} % La largeur de la ligne verticale entre les colonnes
    \begin{enumerate}
        \item $f(x) = \dfrac{\sqrt{25 - x^2}}{x - 1}$\vspace{0.09cm}
        \item $f(x) = \sqrt{x^2 - x + 2}$ 
        \item $f(x) = \dfrac{\sqrt{x^2 - x - 6}}{\sqrt{x^2 - 3x + 2}}$ 
        \item $f(x) = \begin{cases} 
        -x + \sqrt{x^2 - x} & \text{si } x < 0 \\ 
        \dfrac{x^3 - 5x}{x^2 + 3} & \text{si } x \ge 0 
        \end{cases}$
    \end{enumerate}
\end{multicols}
  \item \textbf{Soit $f(x) = \dfrac{4x^2 + 1}{2x^2 + 1}$ et $g(x) = 3x - 2$.}\hfill \textbf{1pt}
\begin{multicols}{2}
\setlength{\columnseprule}{0.9mm} % La largeur de la ligne verticale entre les colonnes
\begin{enumerate}   
    \item Calculer $(g \circ f)(x)$.
    \item Calculer $(f \circ g)(x)$.
\end{enumerate}
\end{multicols}
\item Soit $f$ telle que $f(x) = |x + 2| + |x - 2| + 2x$.

 \begin{enumerate} 
    \item Exprimer $f(x)$ sans le symbole de la valeur absolue 
    
    suivant les valeurs de $x$. \hfill \textbf{0,75pt} 
    \item Déterminer la restriction de $f$ à $]-\infty ; -2]$. \hfill \textbf{0,5pt} \\
\end{enumerate}

\end{enumerate}
\vspace{3cm}
\section*{\underline{ Exercice 3 :}  (7,5 pts) \small\itshape Équations, inéquations et systèmes : méthodes algébriques}

\begin{enumerate}
\item Résoudre dans $\mathbb{R}$ les équations et inéquations suivantes : \textbf{(0,75pt $\times$ 4)}
\begin{multicols}{2}
\setlength{\columnseprule}{0.9mm} % La largeur de la ligne verticale entre les colonnes
\begin{enumerate}
    \item $\sqrt{2x - 3} = x + 2$
    \item $2x + \sqrt{x^2 - 4x - 5} = 18$
    \item $\sqrt{x - 2} \le x - 3$
    \item $\sqrt{x^2 - x + 1} \ge 2 - x$
\end{enumerate}
\end{multicols}
    \item Déterminer le polynôme du second degré $f(x)$ tel que : \hfill \hfill \textbf{(1,5pt)} \\
    $f(1) = 0 \quad ; \quad f(-2) = -3 \quad \text{et} \quad f(3) = 22$
    \item 
    \begin{enumerate}
        \item Résoudre dans $\mathbb{R}^3$ en utilisant la méthode du pivot de Gauss le système : \hfill \textbf{(1,5pt)} \\
        $\begin{cases} 
        x + 2y + z = 11 \\ 
        2x + y - z = 1 \\ 
        4x - 3y + z = 7 
        \end{cases} $

        \item En déduire la résolution des systèmes :\hfill \hfill \textbf{(1,5pt)} \\
        \begin{center}
        (S) $\begin{cases} 
        \dfrac{1}{x+3} + \dfrac{2}{y+2} + \dfrac{1}{z} = 11 \\[10pt] 
        \dfrac{2}{x+3} + \dfrac{1}{y+2} - \dfrac{1}{z} = 1 \\[10pt] 
        \dfrac{4}{x+3} - \dfrac{3}{y+2} + \dfrac{1}{z} = 7 
        \end{cases}$ 
        \quad et \quad 
        (S') $\begin{cases} 
        |x| + 2y^2 + \dfrac{1}{z} = 11 \\[10pt] 
        2|x| + y^2 - \dfrac{1}{z} = 1 \\[10pt] 
        4|x| - 3y^2 + \dfrac{1}{z} = 7 
        \end{cases} $
        \end{center}
    \end{enumerate}
\end{enumerate}

\section*{\underline{ Exercice 4 :} (3 pts) \small\itshape Équations du second degré à paramètre et discussion}

\begin{enumerate}
    \item \textbf{Soit l'équation $(E) : x^2 - (2m + 1)x + m^2 - 2 = 0$}
    
    \begin{enumerate}
        \item Discuter selon les valeurs de $m$, les solutions de l'équation $(E)$ \hfill \textbf{(1,5 pt)}
        
        \item Dans le cas où $(E)$ admet deux solutions $x'$ et $x''$, déterminer $m$ tel que :
        \begin{enumerate}
            \item $x'^2 + x''^2 = 35$ \hfill \textbf{(0,75 pt)}
            \item $|x' - x''| = 5$ \hfill \textbf{(0,75 pt)}
        \end{enumerate}
    \end{enumerate}

\end{enumerate}

\end{document}
